% Revised project proposal (Group 6) -- addresses the feedback of 8 Oct 2026:
% (1) entity = edition, (2) complete DBpedia download with dbp: properties,
% (3) documented Amazon selection with random baseline.
\documentclass[runningheads]{llncs}
\usepackage[T1]{fontenc}
\usepackage{graphicx}
\usepackage{tabularx}
\usepackage{xurl}
\usepackage{booktabs}
\usepackage{array}
\usepackage{threeparttable}
\usepackage{ragged2e}
\setlength{\textfloatsep}{8pt plus 2pt minus 2pt}
\setlength{\intextsep}{8pt plus 2pt minus 2pt}
\setlength{\floatsep}{6pt plus 2pt minus 2pt}
\setlength{\emergencystretch}{2em}

\begin{document}
\title{Web Data Integration Project Proposal (Group 6) -- Revised}
\titlerunning{Book Data Integration}
\authorrunning{Bolduan et al.}
\author{Lucia Bolduan \orcidID{2261625} \and Safa Shaghely \orcidID{2480414} \and Selma Nezihoglu \orcidID{2390917} \and Minji Oh \orcidID{2389125} \and Kejsi Limaci \orcidID{2389812} \and Shiqi Xu \orcidID{1863435}}
\institute{University of Mannheim}
\maketitle

\begin{abstract}
We integrate book data from Goodreads (52,478 records), DBpedia (59,270) and
Amazon (40,000) into a common schema. Following the feedback, the entity is a
book \emph{edition} (e.g.\ the 2018 Putnam hardcover of \emph{Where the Crawdads
Sing}), not a work. Using the ISBN only as an evaluation key, 6,087 editions
occur in at least two datasets. Many works appear in several editions with the
same title and author (6,588 works with 33,540 editions in the Amazon subset),
which makes telling editions apart the central challenge for identity
resolution. The integrated schema has 14 attributes, 12 of which occur in at
least two datasets; \texttt{authors} and \texttt{genres} are list attributes.
\end{abstract}

\section{Use Case and Entity Definition}

The goal is to integrate book information from three independently maintained
sources that describe overlapping books with different attribute names,
formats and levels of detail. One record of the integrated dataset describes
one \textbf{edition}: a specific publication of a work by a publisher, in a
format (hardcover, paperback, \dots), at a date, with its own page count and ISBN.
Goodreads and Amazon list editions directly. A DBpedia resource describes a
work, but its infobox values (ISBN, publisher, release date, pages) are those
of the first edition, so each DBpedia record is treated as the first edition of
its work. Matching therefore has to distinguish, e.g., the hardcover, paperback,
large-print and movie tie-in editions of the same novel, which share title and
author but differ in publisher, date, format and page count.

\section{Datasets}

\textbf{Goodreads.} The ``Best Books Ever'' dataset (scraped from the Goodreads
list of the same name) lists one edition per book. We use the original CSV;
an earlier copy had been re-saved with a spreadsheet program, which had turned
91\% of the ISBNs into scientific notation and altered dates.

\textbf{DBpedia.} All 59,270 resources of type \texttt{dbo:Book} were
downloaded from the public SPARQL endpoint. Because the endpoint silently cuts
off long-running queries, the download (1) lists the book URIs in partitions by
first character, paged with \texttt{LIMIT 10000 OFFSET} (below the 40,000 offset
limit), and (2) fetches the attributes for 150 known URIs per query
(\texttt{VALUES}); a block that returns a partial result (HTTP 206) is split
and fetched again. The final count equals \texttt{COUNT(dbo:Book)}. Year and
language combine \texttt{dbo:releaseDate}/\texttt{dbo:language} with the raw
infobox properties \texttt{dbp:releaseDate}, \texttt{dbp:pubDate} and
\texttt{dbp:language}, which lowers missing values from 88.5\% to 11.4\%
(year) and from 87.6\% to 20.4\% (language). Unlike the first version, books
without an ISBN are no longer excluded, so 39.0\% of the records have no ISBN.

\textbf{Amazon.} The metadata file of Amazon Reviews 2023 (McAuley Lab)
contains 4,448,181 book products. 4,336,308 have a title and an author, and
3,526,297 of these also have an ISBN; this is the population we sample from.
Products without ISBN (mostly Kindle and Audible) are excluded because the
ISBN is our gold-standard key. From this population we take the
\textbf{40,000 editions with the most ratings} (at least 6,506 ratings each).
The selection uses only Amazon's own rating counts, no information from the
other datasets. We chose popularity because Goodreads ``Best Books Ever'' is
itself a popularity-based list: a uniform random sample of 40,000 editions
(fixed seed) shares only 331 editions with Goodreads, compared with 2,169 for
the popularity-based subset (Table~\ref{tab:overlap}). Amazon reports one
rating count for all editions of a work, so the selection contains popular
works together with many of their editions -- exactly the hard cases for
edition-level matching.

\begin{table}[htbp]
\centering
\begin{threeparttable}
\caption{Overview of the datasets.}
\label{tab:overview}
\footnotesize
\newcolumntype{L}[1]{>{\RaggedRight}p{#1}}
\newcolumntype{R}[1]{>{\raggedleft}p{#1}}
\newcolumntype{Y}{>{\RaggedRight\hbadness=10000}X}
\begin{tabularx}{\textwidth}{|L{1.7cm}|L{1.6cm}|l|R{1.0cm}|R{0.8cm}|Y|}
\hline
\textbf{Dataset} & \textbf{Source} & \textbf{Format} & \textbf{\# Ent.} & \textbf{\# Attr.} & \textbf{List of attributes} \\
\hline
Goodreads Best Books Ever & Kaggle / Zenodo\tnote{1} & CSV & 52,478 & 25 & bookId, title, series (MV), author, rating, description, language, isbn, genres, characters (MV), bookFormat, edition (MV), pages, publisher, publishDate, firstPublishDate (MV), awards (MV), numRatings, ratingsByStars, likedPercent, setting (MV), coverImg, bbeScore, bbeVotes, price \\
\hline
DBpedia & SPARQL endpoint\tnote{2} & XML & 59,270 & 9\tnote{4} & title, authors, publisher, release\_date, year, pages (MV), language, genres (MV), isbns (MV) \\
\hline
Amazon Reviews 2023 & Hugging Face\tnote{3} & CSV & 40,000 & 14\tnote{4} & title, authors, format, publisher, edition (MV), publication\_date, language, pages, genres, average\_rating, ratings\_count, price, isbn13, isbn10 \\
\hline
\end{tabularx}
\begin{tablenotes}\footnotesize
  \item[1] \url{https://www.kaggle.com/datasets/arnabchaki/goodreads-best-books-ever}, original: \url{https://doi.org/10.5281/zenodo.4265096}
  \item[2] \url{https://dbpedia.org/sparql}, script \texttt{dbpedia\_books.py}
  \item[3] \url{https://huggingface.co/datasets/McAuley-Lab/Amazon-Reviews-2023}, script \texttt{amazon\_subset.py}
  \item[4] Without the identifiers \texttt{uri} (DBpedia) and \texttt{amazon\_id}, \texttt{asin} (Amazon).
\end{tablenotes}
\end{threeparttable}
\end{table}

Attributes with more than 30\% missing values are marked (MV). In Goodreads
these are \texttt{series} (55.3\%), \texttt{characters} (73.8\%),
\texttt{edition} (90.6\%), \texttt{firstPublishDate} (40.6\%),
\texttt{awards} (79.8\%) and \texttt{setting} (77.9\%); in DBpedia
\texttt{pages} (38.4\%), \texttt{genres} (36.4\%) and \texttt{isbns} (39.0\%);
in Amazon only \texttt{edition} (40.1\%).

\section{Integrated Schema}

Table~\ref{tab:schema} maps the source attributes to the integrated schema.
Compared with the first proposal, \texttt{format} and \texttt{edition} were
added because they distinguish editions of the same work, and the publication
date is kept per edition instead of being reduced to a year.

\begin{table}[htbp]
\centering
\begin{threeparttable}
\caption{Attribute intersection with the integrated schema (one record per edition).}
\label{tab:schema}
\footnotesize\setlength{\tabcolsep}{4pt}
\begin{tabularx}{\textwidth}{|l|l|>{\RaggedRight\arraybackslash}X|}
\hline
\textbf{Attribute name} & \textbf{Type} & \textbf{Datasets in which the attribute is found} \\
\hline
title             & string           & GR (title), DB (title), AM (title) \\
authors           & list$<$string$>$ & GR (author), DB (authors), AM (authors) \\
genres            & list$<$string$>$ & GR (genres), DB (genres), AM (genres) \\
publication\_date & date             & GR (publishDate), DB (release\_date), AM (publication\_date) \\
publisher         & string           & GR (publisher), DB (publisher), AM (publisher) \\
page\_count       & integer          & GR (pages), DB (pages), AM (pages) \\
language          & string           & GR (language), DB (language), AM (language) \\
format            & string           & GR (bookFormat), AM (format) \\
edition           & string           & GR (edition), AM (edition) \\
rating            & float            & GR (rating), AM (average\_rating) \\
ratings\_count    & integer          & GR (numRatings), AM (ratings\_count) \\
price             & float            & GR (price), AM (price) \\
series            & string           & GR (series) \\
awards            & list$<$string$>$ & GR (awards) \\
\hline
\multicolumn{3}{|l|}{\emph{Identifier for overlap estimation and the gold standard only (not used for matching):}} \\
isbn              & string           & GR (isbn), DB (isbns), AM (isbn13, isbn10) \\
\hline
\end{tabularx}
\begin{tablenotes}[para,flushleft]
  \footnotesize GR = Goodreads, DB = DBpedia, AM = Amazon; source column names in brackets.
\end{tablenotes}
\end{threeparttable}
\end{table}

Twelve attributes occur in at least two datasets and seven in all three
(\texttt{title}, \texttt{authors}, \texttt{genres}, \texttt{publication\_date},
\texttt{publisher}, \texttt{page\_count}, \texttt{language}), which allows
fusion by voting. Representations differ: Amazon titles carry series or edition
details in brackets (``Where the Crawdads Sing (Movie Tie-in)''), Goodreads
stores contributors and roles in one string (``J.K. Rowling, Mary GrandPr\'e
(Illustrator)''), DBpedia dates range from years to full dates, and the
sources use different genre vocabularies.

\section{Overlap Estimation}

Two records describe the same edition if they share an ISBN (ISBN-10 and
ISBN-13 normalised to ISBN-13). The ISBN is used only for this estimate and as
the gold standard; the matchers in Phase~II will not use it. Since records
without an ISBN cannot be counted, the numbers in Table~\ref{tab:overlap} are
lower bounds.

\begin{table}[htbp]
\centering
\caption{Edition overlap (shared ISBN-13) and comparison with a random Amazon sample.}
\label{tab:overlap}
\footnotesize
\begin{tabular}{lr}
\toprule
\textbf{Measure} & \textbf{Editions} \\
\midrule
Amazon $\cap$ Goodreads & 2,169 \\
Amazon $\cap$ DBpedia & 1,725 \\
Goodreads $\cap$ DBpedia & 3,271 \\
In all three datasets & 539 \\
\textbf{In at least two datasets} & \textbf{6,087} \\
\midrule
Random Amazon sample (40,000, seed 42) $\cap$ Goodreads & 331 \\
\bottomrule
\end{tabular}
\end{table}

With 6,087 editions in at least two datasets, the requirement of 1,000 is met
even under this strict criterion. The total number of distinct editions is
well above 2,500, since the three sources contain 151,748 records with at most
6,087 shared editions.

To show the difficulty of edition-level matching, we also compared works
(normalised title and first-author surname). Goodreads and DBpedia share
10,619 works; for 4,759 of the 7,561 shared works with ISBNs on both sides
(63\%), the two sources list \emph{different} editions. Amazon and Goodreads
share 3,146 works (22\% in different editions), and within the Amazon subset
6,588 works appear in more than one edition (33,540 records). A matcher based
on title and author alone would therefore produce many false matches; date,
publisher, format and page count are needed to tell editions apart.

\end{document}
